\begin{frame}{NSAK as AI Harness}
    \framesubtitle{The framework as the agent's execution environment}

    \begin{columns}[T,onlytextwidth]
        \begin{column}{0.55\textwidth}
            The agent does not run standalone --\\ \textbf{NSAK is the harness around it}:
            \vspace{0.4em}
            \begin{itemize}[<+->]
                \item<1->[\faServer] \textbf{Execution host} \scriptsize Device at a strategic network position, isolated in a container
                \item<2->[\faCogs] \textbf{Drills as tools} \scriptsize every drill auto-wrapped as a type-safe LangChain tool
                \item<3->[\faLayerGroup] \textbf{Scenarios as entrypoints} \scriptsize AI scenarios run interactive, multi and single agent
                \item<4->[\faSlidersH] \textbf{Configuration} \scriptsize provider, model and credentials via runtime config
                \item<5->[\faPowerOff] \textbf{Safety \& observability} \scriptsize kill switch, structured results, Grafana / Loki logging
            \end{itemize}
        \end{column}

        \begin{column}{0.45\textwidth}
            \onslide<6->{
                \centering
                \footnotesize
                \begin{block}{Reconnaissance in the attack chain}
                    \scriptsize
                    Host discovery $\rightarrow$ port scanning $\rightarrow$ service enumeration
                    $\rightarrow$ (credential access / lateral movement)
                \end{block}

                \vspace{0.6em}
                {\scriptsize\textit{The agent reuses NSAK building blocks instead of ad-hoc scripts.}}
            }
        \end{column}
    \end{columns}
\end{frame}
